% K=3 Symmetric Confusion Matrix: Monotonicity Proof
% Theorem statement is in method.tex; this file contains the proof and remarks only.

\begin{proof}
All three parts reduce to a single lemma about the between-group variance
$V_B(\delta)=\operatorname{Var}\bigl(\mathbb{E}[g(A)\mid A^*_\delta]\bigr)$.

\paragraph{Step 1: Structural reduction.}
The symmetric $C_0$ decomposes as $C_0=(1{-}3c)\,I+c\,\mathbf{1}\mathbf{1}^\top$, so
\[
  C(\delta) = (1{-}\delta)I + \delta\bigl[(1{-}3c)\,I+c\,\mathbf{1}\mathbf{1}^\top\bigr]
  = (1{-}3c\delta)\,I + c\delta\,\mathbf{1}\mathbf{1}^\top.
\]
Setting $s=c\delta\in[0,c]$, each entry is
$C(\delta)_{j\ell}=(1{-}3s)\,\mathbf{1}_{j=\ell}+s$,
so $C(\delta)$ stays in the one-parameter symmetric family.
In particular:
\begin{align}
  q_j(s) &= \textstyle\sum_k C(\delta)_{jk}\,p_k = (1{-}3s)\,p_j + s, \label{eq:k3q} \\
  \mu_j(s) &\equiv \mathbb{E}[g(A)\mid A^*{=}j]
  = \bar{g} + \frac{(1{-}3s)\,p_j\,h_j}{q_j(s)}, \label{eq:k3mu}
\end{align}
where $h_j=g_j-\bar{g}$, $\bar{g}=\sum_k p_k g_k$.

\paragraph{Step 2: Closed-form $V_B(s)$.}
The between-group variance admits a clean factorization:
\begin{equation}
  V_B(s) = \sum_j q_j\,(\mu_j-\bar{g})^2
  = (1{-}3s)^2\sum_{j=0}^{2}\frac{p_j^2\,h_j^2}{q_j(s)}.
  \label{eq:k3VB}
\end{equation}
\emph{Verification.}
At $s{=}0$: $q_j{=}p_j$, so $V_B(0)=\sum_j p_j h_j^2=\operatorname{Var}(g)$.
At $s{=}\tfrac{1}{3}$: $V_B(\tfrac{1}{3})=0$.

\paragraph{Step 3: $V_B$ is strictly decreasing on $[0,\tfrac{1}{3})$ (pointwise, not just at endpoints).}
Differentiating~\eqref{eq:k3VB} via the product rule:
\begin{equation}
  \frac{dV_B}{ds}
  = (1{-}3s)\sum_{j=0}^{2}\frac{p_j^2\,h_j^2}{q_j^2}\;\bigl[\underbrace{-3(1{-}3s)\,p_j - 1 - 3s}_{\displaystyle\Lambda_j(s)}\bigr].
  \label{eq:k3dVB}
\end{equation}

\emph{Derivation.}
Write $V_B=(1{-}3s)^2 F(s)$ with $F(s)=\sum_j p_j^2 h_j^2/q_j(s)$.
The product rule gives $V_B'=-6(1{-}3s)F+(1{-}3s)^2 F'$.
Since $q_j'(s)=1{-}3p_j$, we have $F'(s)=\sum_j p_j^2 h_j^2(3p_j{-}1)/q_j^2$.
Combining:
\begin{align*}
  V_B'
  &= (1{-}3s)\Bigl[-6\sum_j \frac{p_j^2 h_j^2}{q_j}+(1{-}3s)\sum_j \frac{p_j^2 h_j^2(3p_j{-}1)}{q_j^2}\Bigr] \\
  &= (1{-}3s)\sum_j \frac{p_j^2 h_j^2}{q_j^2}\bigl[-6q_j+(1{-}3s)(3p_j{-}1)\bigr].
\end{align*}
Expanding the bracket:
\[
  -6q_j+(1{-}3s)(3p_j{-}1)
  = -6(1{-}3s)p_j-6s+3(1{-}3s)p_j-(1{-}3s)
  = -3(1{-}3s)p_j-1-3s
  = \Lambda_j(s).
\]

\emph{Sign analysis.}
For $s\in[0,\tfrac{1}{3})$: $(1{-}3s)>0$, $p_j\geq 0$, so
\[
  \Lambda_j(s) = -3(1{-}3s)\,p_j - (1+3s) < 0
  \quad\text{for all } j,
\]
since $1+3s>0$ alone forces $\Lambda_j<0$.
Every term $p_j^2 h_j^2/q_j^2\geq 0$, and the prefactor $(1{-}3s)>0$.
Hence $dV_B/ds\leq 0$ for all $s\in[0,\tfrac{1}{3})$.

\emph{Strict inequality.}
If $\operatorname{Var}(g){>}0$, at least two categories $j$ satisfy $p_j>0$ and $h_j\neq 0$, contributing strictly positive terms.
Since $\Lambda_j<0$, the sum is strictly negative:
$dV_B/ds<0$ for all $s\in(0,\tfrac{1}{3})$.

\paragraph{Step 4: From $V_B$ to $|B|$ --- confounding (a).}
Under proportional effects $f=\lambda g$ (Proposition~\ref{prop:conf_prop}):
\[
  B(C(\delta)) = \frac{\lambda\,V_W(\delta)}{V_W(\delta)+\sigma_T^2},
  \quad V_W(\delta)=\operatorname{Var}(g)-V_B(c\delta).
\]
Step~3 shows $V_B(s)$ is strictly decreasing in $s=c\delta$, so $V_W(\delta)$ is strictly increasing from $0$ to $\operatorname{Var}(g)$.
Since $x\mapsto |\lambda|\,x/(x+\sigma_T^2)$ is strictly increasing for $x\geq 0$:
\[
  |B(C(\delta))| = |\lambda|\,\frac{V_W(\delta)}{V_W(\delta)+\sigma_T^2}
\]
is strictly non-decreasing.

\paragraph{Step 5: Mediation (b).}
Under proportional effects $f=\lambda g$ with $g(k)=\mathbb{E}[T\mid A{=}k]$, the mediation bias
(§\ref{par:med_monotone}) has the identical form
$B=\lambda\,V_W/(V_W+\sigma_\varepsilon^2)$
with $\sigma_\varepsilon^2=\mathbb{E}[\operatorname{Var}(T\mid A)]$ replacing $\sigma_T^2$.
The same argument yields $|B|$ strictly non-decreasing.

\paragraph{Step 6: Collider (c).}
Under proportional effects $h=\lambda g$ (Proposition~\ref{prop:col_prop}):
\[
  B_{\mathrm{col}}(C(\delta)) = -\frac{\lambda\,V_B(c\delta)}{\operatorname{Var}(T)-V_B(c\delta)},
  \qquad
  |B_{\mathrm{col}}| = |\lambda|\,\frac{V_B(c\delta)}{\operatorname{Var}(T)-V_B(c\delta)}.
\]
Since $x\mapsto x/(a{-}x)$ is strictly increasing for $x\in[0,a)$
and $V_B(s)$ is strictly decreasing (Step~3),
$|B_{\mathrm{col}}(C(\delta))|$ is strictly non-increasing:
from $|B_{\mathrm{Berkson}}|=|\lambda|\,V_g/(\operatorname{Var}(T){-}V_g)$ at $\delta{=}0$
to $0$ at $\delta$ such that $V_B\to 0$.
\end{proof}

\begin{remark}[Comparison with the general-$K$ proof]
\label{rem:k3_vs_general}
The general-$K$ monotonicity proof (§\ref{app:monotonicity_proof}) relies on
convexity of $V_B(\delta)$ plus the endpoint condition $V_B'(1)\leq 0$,
the latter verified only for $K{=}2$ analytically and $K\geq 3$ numerically.
The $K{=}3$ symmetric proof is \emph{stronger}:
it establishes $dV_B/ds<0$ at every interior point via the signed decomposition~\eqref{eq:k3dVB},
bypassing both the convexity argument and the endpoint condition entirely.
The key enabler is that the symmetric $C_0$ keeps $C(\delta)$ within the one-parameter family
$(1{-}3s)I+s\,\mathbf{1}\mathbf{1}^\top$, yielding the factored form~\eqref{eq:k3VB}
in which the $(1{-}3s)^2$ prefactor and the $q_j^{-1}$ terms interact to produce
a uniformly signed derivative.
\end{remark}

\begin{remark}[General (non-proportional) effects]
\label{rem:k3_general}
For non-proportional effects ($f\neq\lambda g$), the confounding bias
$B(C)=\mathbb{E}[\operatorname{Cov}(g,f\mid A^*)]/D(C)$
does not factor through $V_W$ alone,
and exact monotonicity of $|B(C(\delta))|$ can fail even for $K{=}3$ symmetric $C_0$.
Numerical verification ($50{,}000$ random configurations) finds violations in ${\sim}10\%$ of cases,
with maximum step deviation ${\sim}0.004$,
confirming $\varepsilon$-approximate monotonicity consistent with the general-$K$ results
(Table~\ref{tab:conjecture_verify}).
\end{remark}
